% MAS1004 Assignment 2: long report.
%
% You fill this in with your agent, in your repository. AGENTS.md tells the
% agent how. Red text in the PDF is what still has to be replaced.
%
% - Pictures go in figures/. Copy them from results/, which is not in git.
% - Outputs go in outputs/: copy experiments.csv, LOG.md and the outputs of
%   check_submission.py and curate.py count here. \outputfile{outputs/...}
%   prints a file exactly as it is. Every number in this report has to come
%   from such an output.
%
% To make the PDF: compress this folder into a zip, and on Overleaf choose
% New Project, Upload Project. Set Menu, Compiler to XeLaTeX, and Recompile.

\documentclass[11pt]{article}
\usepackage[ko]{mas1004}

\reporttitle{과제 2: 긴 보고서}
\studentname{\fillin{이름}}
\studentid{\fillin{학번}}

\begin{document}
\makeheader

리더보드 닉네임: Yang \fillin{닉네임이 맞는지 확인}.
마지막 릴리스: Track A \texttt{track-a-1} \fillin{마감 전 마지막 Track A 릴리스로 바꾸기},
Track B \texttt{track-b-1} \fillin{마감 전 마지막 Track B 릴리스로 바꾸기}.

% ---------------------------------------------------------------------------
\section*{Track A}

\begin{center}\footnotesize
\begin{tabular}{lp{0.55\linewidth}r}
\toprule
tag & 이전 실행과 달라진 점 & validation \\
\midrule
baseline (seed 0) & 스타터 레시피 그대로, 학습 사진 전체 (44,257장) & 0.4956 \\
baseline (seed 1) & 시드만 1로 바꿈 & 0.4942 \\
clip8000 (seed 0) & 레시피는 그대로, 학습 사진 선택만 바꿈 (36,886장, 아래 ``사진을 보고 내린 결정'') & 0.4917 \\
\bottomrule
\end{tabular}
\end{center}

\textbf{baseline.} 스타터 레시피를 시드 0과 1로 두 번 학습했다. 두 번의 검증 점수 차이는
0.0014였고, 도시별로는 타이베이가 0.559와 0.524로 가장 크게(0.035) 달랐다. 시드만 바꿔도
이만큼 움직이므로, 이보다 작은 차이는 우연일 수 있다고 보고 이후 변경을 비교했다.
\fillin{baseline을 두 번 돌린 이유와, 두 점수를 보고 든 생각을 직접}

\textbf{clip8000.} 학습 사진에서 지역을 알려주는 단서가 없는 사진(글씨나 건물이 보이지
않는 도로, 사물만 또는 인물만 있는 사진, 어디서나 볼 수 있는 음식)을 7,371장 뺐다.
검증 점수는 0.4917로 baseline 두 번(0.4956, 0.4942)보다 0.003$\sim$0.004 낮았다.
도시별로는 오사카가 0.378로 baseline(0.345, 0.371)보다 높았고, 타이베이는 0.504로
baseline(0.559, 0.524)보다 낮았다. 타이베이는 학습 사진의 25\%가 빠져 다섯 도시 중 가장
많이 빠진 도시다. \fillin{이 결과가 무엇을 뜻한다고 생각하는지 직접}

\fillin{마지막 Track A 레시피와 선택, 그 리더보드 점수를 검증 점수 옆에.
현재 리더보드의 스타터 레시피(track-a-1)는 0.407이고, 채점 서버가 잰 검증 점수 평균은 0.498이다.}

% ---------------------------------------------------------------------------
\section*{사진을 보고 내린 결정}

\picturefile{figures/mistakes_baseline_seed0_busan_00.png}

\textbf{본 것.} baseline 시드 0 모델이 가장 자신 있게 틀린 검증 사진을 도시별로 한 장씩
(\texttt{curate.py mistakes}) 보고, 틀린 사진의 공통점을 적었다.
\begin{itemize}
\item 틀린 사진 중에는 도로나 길을 찍은 사진이 많다.
\item 사진 속 글자가 잘 보이지 않는다.
\item 자동차나 비행기 같은 사물이 많이 찍혀 있다.
\item 특색이 없어 어디서나 볼 수 있는 음식 사진이 섞여 있다.
\end{itemize}

\textbf{규칙.} 행동하기 전에 규칙을 정했다. ``글씨나 건물이 보이지 않는 도로 사진은 뺀다.
사물만 있거나 인물만 있는 사진도 뺀다.'' 그 뒤 음식에 대한 규칙을 더했다. ``음식 사진 중
카페나 브런치처럼 어디서나 볼 수 있는 음식은 빼고, 그 지역의 전통 음식이나 유명한 음식은
남긴다.''

\textbf{적용 1: 직접 고르기.} 부산 학습 사진을 50장마다 한 장씩 뽑은 시트 3장(104장)에서
규칙에 맞는 사진 20장을 직접 골라 뺐다.

\picturefile{figures/look_busan_01.png}

\textbf{적용 2: CLIP으로 넓히기.} 손으로는 학습 사진 44,257장 중 일부밖에 볼 수 없어서,
규칙의 적용은 CLIP(openai/clip-vit-base-patch32)에 맡기고 기준점은 직접 정했다.
CLIP에는 규칙을 문장으로 주었다. 뺄 쪽은 ``a photo of a plain road with no signs and no
buildings'', ``a close-up photo of a single object'', ``a photo of only people'', 남길 쪽은
``a photo of a city street with shop signs and writing'', ``a photo of buildings in a city''이다.
사진마다 뺄 쪽 문장에 가까운 정도를 점수로 매기고(\texttt{scripts/clip\_filter.py}),
순위 구간별로 48장씩 뽑은 시트를 2,000위부터 9,000위까지 500위 단위로 본 뒤
8,000위까지 빼기로 정했다. CLIP은 학습 사진 목록을 만드는 데만 썼고, 학습에는 들어가지 않는다.

\sidebyside{figures/clip_rank_1-250_00.png}{CLIP 순위 1$\sim$250위에서 뽑은 48장}%
{figures/clip_rank_7501-8000_00.png}{CLIP 순위 7,501$\sim$8,000위에서 뽑은 48장}

8,000위 안의 음식 사진은 CLIP으로 한 번 더 나누었다. 음식으로 판단된 833장을
특색 없는 음식(커피, 카페 디저트, 빵, 패스트푸드)과 지역 음식(한국 전통 음식, 부산
돼지국밥·밀면, 일본 라멘·스시, 오사카 오코노미야키·타코야키, 대만 샤오롱바오·야시장 음식)
문장에 견주어 지역 음식 점수를 매겼다. 시트를 보고 0.75 이상이거나 Wikimedia 제목에 음식
이름이 있는 649장은 남기고 184장은 뺐다. \fillin{8,000위와 0.75를 기준으로 고른 이유를 직접}

\sidebyside{figures/food_local_00.png}{지역 음식 점수가 가장 높은 48장 (남김)}%
{figures/food_plain_00.png}{지역 음식 점수가 가장 낮은 48장 (뺌)}

\textbf{결과.} 학습 사진은 44,257장에서 36,886장이 되었다. 검증 점수는 0.4956(baseline
시드 0)에서 0.4917이 되었다.

\outputfile{outputs/curate_count.txt}

\textbf{검증 사진에는 이런 사진이 얼마나 있나.} 같은 CLIP 점수를 검증 사진에도 매겨,
학습 사진에서 고른 기준점(8,000위의 점수 0.9862) 이상인 사진의 비율을 비교했다
(\texttt{scripts/clip\_compare.py}). 검증 사진은 고르는 데 쓰지 않고 세기만 했다.
학습 사진은 18.1\%, 검증 사진은 11.1\%가 기준점 이상이었다. 타이베이는 학습 27.1\%,
검증 16.6\%로 두 쪽 모두 가장 높았다.

\outputfile{outputs/clip_compare.txt}

\fillin{이 비교를 보고 clip8000의 검증 점수를 어떻게 해석하는지 직접}

\textbf{에이전트에게 한 요청.} 이 결정을 위해 에이전트에게 한 요청을 그대로 옮긴다.
\begin{itemize}
\item ``도로,길 사진이 많음, 글자가 잘 보이지 않음, 자동차나 비행기 등의 사물이 많음,
  특색이 없는 어디서나 볼 수 있는 음식 사진이 있음 -> 말로 풀어서 작성해줘''
\item ``글씨나 건물이 보이지 않는 도로 사진 빼기''
\item ``사물만 인물만 있는 사진도 규칙에 추가'' 그리고 시트별 번호 (look\_busan\_00:
  13 14 17 35, look\_busan\_01: 1 4 12 13 15 17 23 24 30 33 34 35, look\_busan\_02: 1 2 5 6)
\item ``너가 알아서 빼줄 수는 없는거야?'' $\rightarrow$ 에이전트가 직접 고르는 것은 과제
  규칙에 맞지 않아, 규칙은 내가 정하고 적용은 CLIP으로 하는 방법을 골랐다.
\item ``이미 손으로 한 건 적용하고 나머지를 진행해''
\item ``음식 사진 중에서 카페 같은 특색 없는 사진을 빼고 지역의 전통음식, 유명한 음식은
  남기고 싶은데 어떻게 해야 하지?'', ``진행해줘''
\item ``0.75 이상'' (지역 음식 기준), ``8000위까지 빼자'' (CLIP 순위 기준)
\end{itemize}

% ---------------------------------------------------------------------------
\section*{Track B}

\textbf{b1.} 스타터 모델 \texttt{train\_b.py}를 그대로 학습했다. ResNet18, 224px, 4에폭,
학습 사진 전체. 검증 점수는 0.4832다 (seoul 0.564, busan 0.369, tokyo 0.604, osaka 0.354,
taipei 0.524).

\fillin{그다음 시도한 변경과, 변경마다의 검증 점수}

\fillin{직접 모은 사진이 있다면 어떻게 모았는지}

\textbf{검증 점수와 리더보드 점수.} b1의 검증 점수는 0.483, 리더보드 점수는 0.397이다.
같은 코드를 채점 서버가 돌린 ``train\_b.py as given''은 리더보드 0.403이었다. 리더보드 사진은
스타터 데이터에 없는 유튜브 채널의 거리 영상 프레임이고, 검증 사진은 스타터 데이터와 같은
출처(Wikimedia, OSV-5M, 유튜브)에서 왔다. \fillin{두 점수가 왜 다르다고 생각하는지 직접}

\subsection*{채점 서버가 보는 제출 파일}
\outputfile{outputs/check_submission.txt}

% ---------------------------------------------------------------------------
\section*{무엇이 차이를 만들었나}

\fillin{두 트랙에서 가장 도움이 된 변경과 가장 덜 된 변경, 숫자와 함께}

\fillin{일주일이 더 있다면 무엇을 하고 싶은지}

% ---------------------------------------------------------------------------
\section*{모든 실행}
\outputfile{outputs/experiments.csv}

\section*{실험 노트}
\outputfile{outputs/LOG.md}

\end{document}
